
\chapter{버전별 변경사항}

\newversionsection{0.5.2}{2025-11-28}

\begin{itemize}
    \item \textbf{추가된 기능}
    \begin{itemize}
        \item 수정된 ASHRAE 140 케이스 600 모델을 예시파일로 추가.
    \end{itemize}
    
    \item \textbf{변경된 기능}
    \par -

    \item \textbf{수정된 기능}
    \par -
    
    \item \textbf{제거된 기능}
    \par -
\end{itemize}


\newversionsection{0.5.1}{2025-11-17}
*\textit{V0.5.0을 기준으로 작성되었습니다.}

\begin{itemize}
    \item \textbf{추가된 기능}
    \begin{itemize}
        \item 공급설비의 종류로 전기바닥난방(electric\_radiant\_floor)을 추가 (electric\_radiator와 동일하게, 이름 외 별도 속성 변수 없음).
    \end{itemize}
    
    \item \textbf{변경된 기능}
    \begin{itemize}
        \item 냉난방 설비의 자동 용량 산정 (Autosizing) 기간을 1월 15일부터 8월 15일까지에서 1월 한달간 및 8월 한달간으로 변경. 이로 인해 용량 미입력시 자동 계산되던 설비의 용량이 변화할 수 있음.
        \item 패키지에어컨이 EnergyPlus의 AirConditioner:VariableRefrigerantFlow 객체를 사용 (기존에는 ZoneHVAC:WindowAirConditioner사용; 공조기+EHP와 동일한 객체 사용하도록 변경).
        \item 패키지에어컨 및 EHP 사용시 냉방 코일의 현열비 및 정격 풍량을 자동 산정하지 않고 각각 0.7, 10L/$s\cdot m^2$으로 고정.
    \end{itemize}

    \item \textbf{수정된 기능}
    \begin{itemize}
        \item 패키지에어컨 사용시 냉방에너지가 계산되지 않던 (EnergyPlus 시뮬레이션에서 항상 가동하지 않던) 오류 수정.
        \item EHP 사용시 특정 조건에서 발생하던 EnergyPlus 실행 오류 수정.
        \item 난방유 사용시 grr 파일에서 결과에 집계되지 않던 오류 수정.
        \item EnergyPlus 실행 후 err 파일을 읽을 수 없던 경우 발생하던 오류를 EnergyPlusError Exception으로 예외처리 가능하도록 수정.
    \end{itemize}
    
    \item \textbf{제거된 기능}
    \par -
\end{itemize}